\documentclass[10pt,a4paper]{amsart}
\usepackage[margin=19mm]{geometry}
\usepackage{amsmath,amssymb,mathtools}
\usepackage[scr=rsfs,cal=euler]{mathalfa}
\usepackage{xfrac}
\usepackage[unicode,hidelinks,bookmarks=false]{hyperref}
\newcommand{\ie}{i.e.\ }
\newcommand{\cf}{cf.\ }
\newcommand{\resp}{respectively\ }
\newcommand{\Subsection}{\subsection}
\providecommand{\nobreakdash}{\nobreak\hbox{-}}
\setlength{\emergencystretch}{2em}
\multlinegap0pt
\usepackage{amssymb}
\usepackage{tikz}
\usepackage{booktabs}
\newtheorem{theorem}{Theorem}
\newcommand{\C}{{\mathbb C}}
\newcommand{\PP}{{\mathbb P}}
\title{Unbounded Reinhardt domains with finite-dimensional Bergman spaces in $\C^n$}
\author{Chika Hayashida, Joe Kamimoto}
\date{Source: arXiv:2602.13732v1 (CC BY 4.0)}
\begin{document}
\maketitle

\noindent\emph{Source and licence.} Unbounded Reinhardt domains with finite-dimensional Bergman spaces in $\C^n$. Version arXiv:2602.13732v1 (CC BY 4.0), distributed by arXiv under the Creative Commons Attribution 4.0 International licence, http://creativecommons.org/licenses/by/4.0/, as recorded in the arXiv licence metadata for this exact version and shown on the abstract page https://arxiv.org/abs/2602.13732v1. Full licence notice: \texttt{licenses/arxiv-2602.13732-CC-BY-4.0.txt}.\par\smallskip

\noindent\textit{Editorial scope.} Counted unit: the article's theorem theorem (source0.tex lines 568--582). The article's complete original proof of it is delivered with it (source0.tex lines 585--633, one \texttt{proof} environment of 49 lines). No mathematics is written, repaired or re-derived: the statement and the proof are the article's own lines, in the article's own notation, and no retained line is substituted or re-flowed.  No further source-local context is needed: every cross-reference of every retained line resolves inside the two delivered spans.  Nothing was omitted from any retained span, so the package carries no omission record.  No retained line contains a citation, so the wrapper carries no bibliography and no bibliography entry.  No retained line contains a figure environment, an image-inclusion command or drawing code, so no image is delivered and the package declares no asset dependency.  Editorial formatting only, in addition: an AMS \texttt{amsart} preamble that restates the article's own declarations and macros used by the retained lines, namely the environments \texttt{theorem} on separate counters, together with the standard packages those lines need.  The retained environments print in this document as Theorem 1 in order of appearance, whereas the article prints the same items with the numbering of its own full document.  The complete native source of the article as distributed in the arXiv e-print for this exact version is bundled unchanged as \texttt{sources/194661/source0.tex}.
\begin{theorem}
If 
%\begin{equation}\label{eqn:3.1}
$\frac{2}{n-1}<a\leq\frac{3}{n-1}$,
%\end{equation} 
then the following statements hold: 
\begin{enumerate}
\item
The Bergman space is 
$A^2(D(a))={\rm Span}\{1,z_1,\ldots,z_n\}.$
\item
The Bergman metric of $D(a)$ has constant 
holomorphic sectional curvature equal to $2$.
\end{enumerate}
\end{theorem}

\begin{proof}

(i) \;\; It follows from Theorem~2.1 that 
the condition 
$1<(n-1)a-1\leq 2$ is equivalent to (i). 

(ii) \;\;
From (i), 
the Bergman kernel of $D(a)$ can be written
as 
\begin{equation}
K_{D(a)}(z)=c_0+\sum_{k=1}^n c_k |z_k|^2
\quad (z\in D(a)), 
\end{equation}
where 
$$
c_0={\rm vol} (D(a))^{-1}, 
\quad 
c_k=\|z_k\|^{-2}=
\left(\int_{D(a)} |z_k|^2 dV(z)\right)^{-1}.$$ 
%are positive constants depending only on $a$ and $n$.
Furthermore, we have
$$
K_{D(a)}(z)=
c_0\left(1+\sum_{k=1}^n \tilde{c}_k |z_k|^2\right)
\quad (z\in D(a)),
$$ 
where $\tilde{c}_k=c_k/c_0$ for
$k=1,\ldots,n$. 
Through
the scaling $(w_1,\ldots,w_n)=
(\sqrt{\tilde{c}_1}z_1,\ldots,
\sqrt{\tilde{c}_n}z_n)$,
the Bergman metric of $D(a)$ 
is isometric to the Fubini-Study metric
of $\PP^n$ restricted to 
a domain in the cell $\C^n$ with 
$$
\omega_{FS} =
\partial\bar{\partial} \log 
\left(
1+\sum_{k=1}^n |w_k|^2
\right). 
$$
Therefore, the Bergman metric of $D(a)$ 
has a constant holomorphic 
sectional curvature equal 
to $2$.
\end{proof}

\end{document}
